\subsection{Orthogonal Classification over Common Fields}

Throughout this subsection, orthogonal geometries are finite-dimensional over a field $\mathbb{F}$; additional hypotheses are stated locally.

\begin{theorem} \label{thm:orthogonal-bases-and-diagonal-congruence} \hlt{Orthogonal Bases and Diagonal Congruence}\\
Let $(\mathcal{V},B)$ be an orthogonal geometry, $q(\mathbf{x})\defeq B(\mathbf{x},\mathbf{x})$, $n=\dim\mathcal{V}$, and $r=\operatorname{rank}B$.
\begin{enumerate}[label=(\roman*)]
\setlength{\itemsep}{0pt}
\item An orthogonal basis $\beta$ with $M_{\beta}=\operatorname{diag}(a_1,\ldots,a_r,0_{n-r})$ and $a_1,\ldots,a_r\in\mathbb{F}^{\times}$ exists $\Longleftrightarrow\operatorname{char}(\mathbb{F})\neq2\ \vee\ B=0\ \vee\ q\not\equiv0$.
\item A symmetric $A\in\operatorname{Mat}(n,n,\mathbb{F})$ is congruent to a diagonal matrix $\Longleftrightarrow\operatorname{char}(\mathbb{F})\neq2\ \vee\ A=0\ \vee\ (\exists\,\mathbf{x}\in\mathbb{F}^n)\,\mathbf{x}^tA\mathbf{x}\neq0$.
\end{enumerate}
\begin{proof}
\begin{enumerate}[label=(\roman*)]
\setlength{\itemsep}{0pt}
\item[$(i)$] ($\Rightarrow$) Only $\operatorname{char}(\mathbb{F})=2$ and $q\equiv0$ require proof. In this case every diagonal entry of $M_{\beta}$ equals $q(\mathbf{v})=0$, so $M_{\beta}=0$ and $B=0$.
\item[$(i)$] ($\Leftarrow$) If $B=0$, any basis works. Assume $B\neq0$. Then $q\not\equiv0$ by hypothesis if $\operatorname{char}(\mathbb{F})=2$, and by Proposition \ref{prop:form-type-relations-by-characteristic}(ii) otherwise.\\
Splitting successive nonisotropic lines by Theorem \ref{thm:metric-orthogonal-complement-dimension}(ii) produces $\mathcal{V}=\mathbb{F}\mathbf{u}_1\odot\cdots\odot\mathbb{F}\mathbf{u}_s\odot\mathcal{W}$ with $q(\mathbf{u}_i)=a_i\in\mathbb{F}^{\times}$ and $q|_{\mathcal{W}}\equiv0$.\\
If $\operatorname{char}(\mathbb{F})\neq2$, then $B|_{\mathcal{W}}=0$ by Proposition \ref{prop:form-type-relations-by-characteristic}(ii). Appending any basis of $\mathcal{W}$ produces a diagonal Gram matrix of rank $r$.\\
Assume $\operatorname{char}(\mathbb{F})=2$; then $B|_{\mathcal{W}}$ alternating. By Theorem \ref{thm:symplectic-geometries-hyperbolic-classification}(ii), $\mathcal{W}=R\odot\mathcal{H}_1\odot\cdots\odot\mathcal{H}_t$, $B|_R=0$, each $\mathcal{H}_j$ hyperbolic.\\
Fix a current $\mathbf{u}$ with $q(\mathbf{u})=a\neq0$. For a hyperbolic basis $(\mathbf{e},\mathbf{f})$ of $\mathcal{H}_j$, rescale $\mathbf{f}$ so that $B(\mathbf{e},\mathbf{f})=a$.\\
The vectors $\mathbf{u}+\mathbf{e}$, $\mathbf{u}+\mathbf{f}$, and $\mathbf{u}+\mathbf{e}+\mathbf{f}$ are independent by coefficient comparison.\\
These vectors have norm $a$ and mutual pairings $a+a=0$, so they form an orthogonal basis of $\mathbb{F}\mathbf{u}\odot\mathcal{H}_j$.\\
Retain one new line for the next plane; after $t$ replacements, append a basis of $R$.\\
The resulting diagonal matrix has exactly $r$ nonzero entries by rank invariance.
\item[$(ii)$] ($\Rightarrow$) Let $D=P^tAP$ be diagonal with $P$ invertible.\\
If $\operatorname{char}(\mathbb{F})=2$ and $(\forall\,\mathbf{x}\in\mathbb{F}^n)\,\mathbf{x}^tA\mathbf{x}=0$, then $(\forall\,j\in\{1,\ldots,n\})\,D_{jj}=(P\mathbf{e}_j)^tA(P\mathbf{e}_j)=0$. Thus $D=0$ and $A=0$.
\item[$(ii)$] ($\Leftarrow$) Define $B_A(\mathbf{x},\mathbf{y})\defeq\mathbf{x}^tA\mathbf{y}$ on $\mathbb{F}^n$. Its norm map is $\mathbf{x}\mapsto\mathbf{x}^tA\mathbf{x}$, so the hypothesis is exactly that of (i).\\
Since $A=A^t$, the form $B_A$ is symmetric, so choose an orthogonal basis by (i).\\
Its change-of-basis matrix diagonalizes $A$ by Proposition \ref{prop:metric-form-matrix-recovery-and-congruence}(ii).
\end{enumerate}
\hlt{Structural Commentary}: Split nonisotropic lines; in characteristic $2$, absorb each residual hyperbolic plane into one retained line.\\
\hlt{Key Insight}: A nonzero alternating symmetric form in characteristic $2$ is the sole obstruction.
\end{proof}
\end{theorem}

\begin{proposition} \label{prop:hyperbolic-plane-orthogonal-criteria} \hlt{Recognizing Orthogonal Hyperbolic Planes}\\
Let $\mathcal{H}$ be a $2$-dimensional nonsingular orthogonal geometry with $\operatorname{char}(\mathbb{F})\neq2$. The following are equivalent:
\begin{enumerate}[label=(\roman*)]
\setlength{\itemsep}{0pt}
\item $\mathcal{H}$ is a hyperbolic plane;
\item $(\exists\,\mathbf{u}\in\mathcal{H}\setminus\{\mathbf{0}\})\,\innerprod{\mathbf{u},\mathbf{u}}=0$;
\item $\Delta(\mathcal{H})=-(\mathbb{F}^{\times})^2$.
\end{enumerate}
\begin{proof}
\begin{enumerate}[label=(\roman*)]
\setlength{\itemsep}{0pt}
\item[$(i)$] $\Rightarrow (ii)$ A hyperbolic basis contains a nonzero isotropic vector by the definition of a hyperbolic pair.
\item[$(ii)$] $\Rightarrow (iii)$ By Theorem \ref{thm:isotropic-vector-sits-in-hyperbolic-plane}, fix $\mathbf{v}$ such that $(\mathbf{u},\mathbf{v})$ is hyperbolic.\\
Since $\dim\mathcal{H}=2$, this pair is a basis. Its Gram matrix has determinant $\det\left(\begin{smallmatrix}0&1\\1&0\end{smallmatrix}\right)=-1$, so $\Delta(\mathcal{H})=-(\mathbb{F}^{\times})^2$.
\item[$(iii)$] $\Rightarrow (i)$ By Theorem \ref{thm:orthogonal-bases-and-diagonal-congruence}(i), write $M=\operatorname{diag}(a,b)$ with $a,b\in\mathbb{F}^{\times}$.\\
Choose $c\in\mathbb{F}^{\times}$ with $ab=-c^2$ by the discriminant hypothesis.\\
Then $a\left(\frac{c}{a}\right)^2+b=0$, so $\left(\frac{c}{a},1\right)$ is a nonzero isotropic coordinate vector.\\
It lies in a hyperbolic plane by Theorem \ref{thm:isotropic-vector-sits-in-hyperbolic-plane}, and that plane equals $\mathcal{H}$ by dimension.
\end{enumerate}
\hlt{Structural Commentary}: Pass from a hyperbolic basis to determinant $-1$, then recover isotropy from the discriminant square class.\\
\hlt{Key Insight}: In a nonsingular orthogonal plane, representing zero is exactly hyperbolicity.
\end{proof}
\end{proposition}

\begin{example} \label{ex:characteristic-two-orthogonal-symplectic-plane} \hlt{Characteristic Two Orthogonal-Symplectic Plane}\\
Assume $\operatorname{char}(\mathbb{F})=2$ and let $(\mathcal{H},B)$ have basis $(\mathbf{e},\mathbf{f})$ with Gram matrix $\left(\begin{smallmatrix}0&1\\1&0\end{smallmatrix}\right)$.\\
The matrix is symmetric and has determinant $1$, while $(\forall\,a,b\in\mathbb{F})\,B(a\mathbf{e}+b\mathbf{f},a\mathbf{e}+b\mathbf{f})=2ab=0$; hence $B$ is nonsingular and alternating.\\
For $\mathbf{0}\neq\mathbf{x}\in\mathcal{H}$, alternation implies $\mathbb{F}\mathbf{x}\leq(\mathbb{F}\mathbf{x})^{\perp}$, and both spaces have dimension $1$ by Theorem \ref{thm:metric-orthogonal-complement-dimension}(i).\\
Thus $(\mathbb{F}\mathbf{x})^{\perp}=\mathbb{F}\mathbf{x}$, so no second basis vector can be orthogonal to $\mathbf{x}$. This is precisely the obstruction in Theorem \ref{thm:orthogonal-bases-and-diagonal-congruence}.
\end{example}

\begin{computationmethod} \label{meth:lagrange-orthogonalization-arbitrary-field} \hlt{Lagrange Orthogonalization}\\
Let $(\mathcal{V},B)$ be an orthogonal geometry with $\operatorname{char}(\mathbb{F})\neq2$ and $q(\mathbf{x})\defeq B(\mathbf{x},\mathbf{x})$.\\
Input: a finite spanning list $L$. Output: an orthogonal basis of $\mathcal{V}$.
\begin{enumerate}[label=(\roman*)]
\setlength{\itemsep}{0pt}
\item Initialize $\mathcal{W}_0=\mathcal{V}$ and $L_0=L$. At stage $k$, stop with $R=\mathcal{W}_k$ if $q|_{\mathcal{W}_k}\equiv0$.\\
If $(\exists\,\mathbf{x}\in L_k)\,q(\mathbf{x})\neq0$, take such an $\mathbf{x}$ as $\mathbf{u}_{k+1}$. Otherwise write $L_k=(\mathbf{x}_1,\ldots,\mathbf{x}_{\ell_k})$.\\
Then $(\exists\,i,j\in\{1,\ldots,\ell_k\})\,B(\mathbf{x}_i,\mathbf{x}_j)\neq0$, since the contrary would imply $B|_{\mathcal{W}_k}=0$ by bilinearity and spanning.\\
Put $\mathbf{u}_{k+1}=\mathbf{x}_i+\mathbf{x}_j$; then $q(\mathbf{u}_{k+1})=2B(\mathbf{x}_i,\mathbf{x}_j)\neq0$.
\item Put $\mathcal{W}_{k+1}=\mathcal{W}_k\cap(\mathbb{F}\mathbf{u}_{k+1})^{\perp}$. Let $L_{k+1}$ be the nonzero vectors among $\left\{\mathbf{x}-\frac{B(\mathbf{x},\mathbf{u}_{k+1})}{q(\mathbf{u}_{k+1})}\mathbf{u}_{k+1}:\mathbf{x}\in L_k\right\}$.\\
By Theorem \ref{thm:metric-orthogonal-complement-dimension}(ii), $\mathcal{W}_k=\mathbb{F}\mathbf{u}_{k+1}\odot\mathcal{W}_{k+1}$. The displayed update is the projection onto $\mathcal{W}_{k+1}$, so its image $L_{k+1}$ spans $\mathcal{W}_{k+1}$.
\item Since $\dim\mathcal{W}_{k+1}=\dim\mathcal{W}_k-1$, the process stops. For any basis $(\mathbf{z}_1,\ldots,\mathbf{z}_m)$ of $R$, output $\beta=(\mathbf{u}_1,\ldots,\mathbf{u}_s,\mathbf{z}_1,\ldots,\mathbf{z}_m)$.\\
Selected lines are mutually orthogonal and orthogonal to $R$. Since $q|_R\equiv0$, $B|_R=0$ by Proposition \ref{prop:form-type-relations-by-characteristic}(ii); hence $\beta$ is orthogonal.
\end{enumerate}
\end{computationmethod}

\begin{definition} \label{def:universal-orthogonal-geometry} \hlt{Universal Orthogonal Geometry}\\
An orthogonal geometry $\mathcal{V}$ over $\mathbb{F}$ is \hlt{universal} if its norm map $q(\mathbf{x})\defeq\innerprod{\mathbf{x},\mathbf{x}}$ satisfies $q(\mathcal{V})=\mathbb{F}$.
\end{definition}

\begin{proposition} \label{prop:finite-field-square-classes-and-universality} \hlt{Finite-Field Square Classes and Universality}\\
Let $\mathbb{F}_q$ be a finite field and put $S=(\mathbb{F}_q^{\times})^2$. For an orthogonal geometry $\mathcal{V}$, write $q_{\mathcal{V}}(\mathbf{x})\defeq\innerprod{\mathbf{x},\mathbf{x}}$.
\begin{enumerate}[label=(\roman*)]
\setlength{\itemsep}{0pt}
\item If $\operatorname{char}(\mathbb{F}_q)=2$, then $x\mapsto x^2$ is a bijection of $\mathbb{F}_q$.
\item If $\operatorname{char}(\mathbb{F}_q)\neq2$, then $|S|=\frac{q-1}{2}$.
\item If $\operatorname{char}(\mathbb{F}_q)\neq2$, then $(\forall\,d\in\mathbb{F}_q^{\times}\setminus S)\,\mathbb{F}_q^{\times}\setminus S=dS$.
\item If $\mathcal{V}$ is orthogonal over odd $\mathbb{F}_q$ and contains a nonsingular subspace of dimension at least $2$, then $q_{\mathcal{V}}(\mathcal{V})=\mathbb{F}_q$.
\end{enumerate}
\begin{proof}
\begin{enumerate}[label=(\roman*)]
\setlength{\itemsep}{0pt}
\item If $x^2=y^2$, then $(x-y)^2=0$ in characteristic $2$, so $x=y$. Thus $x\mapsto x^2$ is injective and hence bijective on the finite set $\mathbb{F}_q$.
\item For $x,y\in\mathbb{F}_q^{\times}$, $x^2=y^2\Longleftrightarrow(x-y)(x+y)=0\Longleftrightarrow x=\pm y$.\\
Since $x\neq-x$, the $\frac{q-1}{2}$ pairs $\{x,-x\}$ are exactly the fibers of $x\mapsto x^2$, so $|S|=\frac{q-1}{2}$.
\item Fix $d\in\mathbb{F}_q^{\times}\setminus S$. If $dS\cap S\neq\varnothing$, choose $s_1,s_2\in S$ with $ds_1=s_2$. Then $d=s_2s_1^{-1}\in S$, a contradiction.\\
Hence $dS\cap S=\varnothing$. Since $|dS|=|S|=\frac{q-1}{2}$ by (ii), $\mathbb{F}_q^{\times}\setminus S=dS$.
\item Put $Q=q_{\mathcal{V}}$. By Theorem \ref{thm:orthogonal-bases-and-diagonal-congruence}(i), the nonsingular subspace contains orthogonal $\mathbf{u},\mathbf{v}$. Put $a\defeq Q(\mathbf{u})\neq0$ and $b\defeq Q(\mathbf{v})\neq0$.\\
Let $T\defeq\{\alpha^2:\alpha\in\mathbb{F}_q\}$. By (ii), $|T|=\frac{q+1}{2}$. For fixed $c\in\mathbb{F}_q$, the sets $aT$ and $c-bT$ both have $\frac{q+1}{2}$ elements, so they intersect.\\
Choose $\alpha,\beta\in\mathbb{F}_q$ with $a\alpha^2=c-b\beta^2$. Orthogonality implies $Q(\alpha\mathbf{u}+\beta\mathbf{v})=a\alpha^2+b\beta^2=c$, so $Q(\mathcal{V})=\mathbb{F}_q$.
\end{enumerate}
\hlt{Structural Commentary}: Count square-map fibers, identify the nonsquare coset, then intersect two affine square-value sets.\\
\hlt{Key Insight}: Two subsets of $\mathbb{F}_q$ of size $\frac{q+1}{2}$ must intersect.
\end{proof}
\end{proposition}

\begin{theorem} \label{thm:orthogonal-canonical-forms-common-fields} \hlt{Canonical Orthogonal Forms over Common Fields}\\
Let $M=M^t\in\operatorname{Mat}(n,n,\mathbb{F})$, let $r=\operatorname{rank}M$, and write $M\sim N$ for congruence. If $\operatorname{char}(\mathbb{F})=2$, assume $M$ is not nonzero alternating.
\begin{enumerate}[label=(\roman*)]
\setlength{\itemsep}{0pt}
\item If $\mathbb{F}$ is algebraically closed, then $M\sim I_r\oplus0_{n-r}$, and rank is a complete congruence invariant.
\item If $\mathbb{F}=\mathbb{R}$, then $M\sim I_p\oplus(-I_m)\oplus0_k$ for unique $p,m,k\in\mathbb{Z}_{\geq0}$ satisfying $p+m+k=n$.
\item If $\mathbb{F}=\mathbb{F}_q$ has characteristic $2$, then $M\sim I_r\oplus0_{n-r}$, and rank is a complete congruence invariant among the admitted matrices.
\item If $\mathbb{F}=\mathbb{F}_q$ has odd characteristic and $d\in\mathbb{F}_q^{\times}$ is a fixed nonsquare, then $r=0$ implies $M\sim0_n$.\\
For $r>0$, exactly one of $M\sim I_r\oplus0_{n-r}$ and $M\sim I_{r-1}\oplus[d]\oplus0_{n-r}$ holds.
\end{enumerate}
\begin{proof}
\begin{enumerate}[label=(\roman*)]
\setlength{\itemsep}{0pt}
\item By Theorem \ref{thm:orthogonal-bases-and-diagonal-congruence}(ii), write $M\sim\operatorname{diag}(a_1,\ldots,a_r,0_{n-r})$.\\
By algebraic closure, choose $(s_1,\ldots,s_r)\in(\mathbb{F}^{\times})^r$ satisfying $(\forall\,i\in\{1,\ldots,r\})\,a_i=s_i^2$.\\
Congruence by $\operatorname{diag}(s_1^{-1},\ldots,s_r^{-1},I_{n-r})$ produces $I_r\oplus0_{n-r}$.\\
Rank is congruence-invariant by Proposition \ref{prop:metric-form-matrix-recovery-and-congruence}(iv), so it is complete.
\item By Theorem \ref{thm:orthogonal-bases-and-diagonal-congruence}(ii), diagonalization and rescaling produce $D=I_p\oplus(-I_m)\oplus0_k$.\\
Let $E_{\leq0}$ be the coordinate subspace on which the form is nonpositive.\\
Every positive-definite subspace intersects $E_{\leq0}$ trivially, so its dimension is at most $\operatorname{codim}E_{\leq0}=p$.\\
The positive coordinate subspace has dimension $p$, so the maximum is $p$ and is congruence-invariant.\\
The same argument for $-D$ determines $m$, and then $k=n-p-m$.
\item By Theorem \ref{thm:orthogonal-bases-and-diagonal-congruence}(ii), write $M\sim\operatorname{diag}(a_1,\ldots,a_r,0_{n-r})$.\\
By Proposition \ref{prop:finite-field-square-classes-and-universality}(i), choose $(s_1,\ldots,s_r)\in(\mathbb{F}_q^{\times})^r$ satisfying $(\forall\,i\in\{1,\ldots,r\})\,a_i=s_i^2$.\\
Congruence by $\operatorname{diag}(s_1^{-1},\ldots,s_r^{-1},I_{n-r})$ produces $I_r\oplus0_{n-r}$, and congruence invariance of rank proves completeness.
\item If $r=0$, then $M=0_n$. Assume $r>0$ and put $S=(\mathbb{F}_q^{\times})^2$.\\
By Theorem \ref{thm:orthogonal-bases-and-diagonal-congruence}(ii), diagonalize $M$ and let $\mathcal{U}$ be the nonsingular block, with norm map $Q$.\\
While $\dim\mathcal{U}\geq2$, choose $\mathbf{u}\in\mathcal{U}$ with $Q(\mathbf{u})=1$ by Proposition \ref{prop:finite-field-square-classes-and-universality}(iv).\\
By Theorem \ref{thm:metric-orthogonal-complement-dimension}(ii),(iv), split $\mathbb{F}\mathbf{u}$ with a nonsingular complement and continue there.\\
The nonsingular block becomes $I_{r-1}\oplus[a]$ with $a\in\mathbb{F}_q^{\times}$. By Proposition \ref{prop:finite-field-square-classes-and-universality}(ii),(iii), exactly one of $a\in S$ and $a\in dS$ holds.\\
Rescaling the final line produces the corresponding stated form.\\
Any congruence maps radicals onto radicals and induces an isometry of the nonsingular quotients.\\
The two quotients have discriminants $S$ and $dS$, which are distinct because $d$ is nonsquare.\\
Discriminant is isometry-invariant by Proposition \ref{prop:metric-form-matrix-recovery-and-congruence}(v), so the normal forms are not congruent.
\end{enumerate}
\hlt{Structural Commentary}: Diagonalize, normalize coefficients by the field, then separate residual forms by inertia or discriminant.\\
\hlt{Key Insight}: Signs over $\mathbb{R}$ and square classes over finite fields survive rescaling.
\end{proof}
\end{theorem}

\begin{definition} \label{def:minkowski-metric-vector-space} \hlt{Minkowski Space}\\
\hlt{Minkowski space} $\mathbb{R}^{3,1}$ is $\mathbb{R}^4$ with standard Gram matrix $\operatorname{diag}(1,1,1,-1)$.
\end{definition}

\begin{example} \label{ex:minkowski-null-cone-and-inertia} \hlt{Null Vectors in Minkowski Space}\\
In $\mathbb{R}^{3,1}$, $\mathcal{N}\defeq\{\mathbf{x}\in\mathbb{R}^4\setminus\{\mathbf{0}\}:x_1^2+x_2^2+x_3^2=x_4^2\}$ is the null cone.\\
The vector $\mathbf{e}_1+\mathbf{e}_4$ lies in $\mathcal{N}$. No isometry to Euclidean $\mathbb{R}^4$ exists: its matrix $A$ would satisfy $A^tA=\operatorname{diag}(1,1,1,-1)$.\\
Taking determinants would force $(\det A)^2=-1$, which is impossible over $\mathbb{R}$.
\end{example}

\begin{example} \label{ex:finite-field-orthogonal-rank-two-classes} \hlt{Two Rank-Two Orthogonal Planes over $\mathbb{F}_5$}\\
In $\mathbb{F}_5$, $(\mathbb{F}_5^{\times})^2=\{1,4\}$ and the nonsquare class is $\{2,3\}$. Let $q_1(x,y)\defeq x^2+y^2$ and $q_2(x,y)\defeq x^2+2y^2$.
\begin{enumerate}[label=(\roman*)]
\setlength{\itemsep}{0pt}
\item The form $q_1$ has matrix $I_2$ and is hyperbolic by $q_1(1,2)=0$ and Proposition \ref{prop:hyperbolic-plane-orthogonal-criteria}.
\item The form $q_2$ is anisotropic: if $q_2(x,y)=0$ and $y=0$, then $x=0$; if $y\neq0$, then $\left(\frac{x}{y}\right)^2=3$, contradicting that $3$ is nonsquare.
\end{enumerate}
Thus the two forms realize the nonsingular rank-$2$ classes in Theorem \ref{thm:orthogonal-canonical-forms-common-fields}(iv).
\end{example}

\begin{lemma} \label{lem:odd-finite-field-anisotropic-residue-bound} \hlt{Odd Finite-Field Isotropy Bound}\\
Let $\mathcal{V}$ be nonsingular orthogonal over odd $\mathbb{F}_q$ and put $q(\mathbf{x})\defeq\innerprod{\mathbf{x},\mathbf{x}}$.\\
Write $\mathcal{H}^{\odot s}$ for an orthogonal sum of $s$ hyperbolic planes.
\begin{enumerate}[label=(\roman*)]
\setlength{\itemsep}{0pt}
\item If $\dim\mathcal{V}\geq3$, then $(\exists\,\mathbf{x}\in\mathcal{V}\setminus\{\mathbf{0}\})\,q(\mathbf{x})=0$.
\item $(\exists\,s\in\mathbb{Z}_{\geq0})(\exists\,\mathcal{R}\leq\mathcal{V})\,\mathcal{V}\cong\mathcal{H}^{\odot s}\odot\mathcal{R}$, where $\mathcal{R}$ is nonsingular anisotropic and $\dim\mathcal{R}\leq2$.
\end{enumerate}
\begin{proof}
\begin{enumerate}[label=(\roman*)]
\setlength{\itemsep}{0pt}
\item By Theorem \ref{thm:orthogonal-bases-and-diagonal-congruence}(i), choose orthogonal nonisotropic vectors $\mathbf{u},\mathbf{v}\in\mathcal{V}$. Put $\mathcal{P}=\operatorname{span}(\mathbf{u},\mathbf{v})$, which is nonsingular.\\
By Theorem \ref{thm:metric-orthogonal-complement-dimension}(i), $\dim\mathcal{P}^{\perp}=\dim\mathcal{V}-2\geq1$; choose $\mathbf{0}\neq\mathbf{z}\in\mathcal{P}^{\perp}$.\\
If $q(\mathbf{z})=0$, take $\mathbf{x}=\mathbf{z}$. Otherwise choose $\mathbf{p}\in\mathcal{P}$ with $q(\mathbf{p})=-q(\mathbf{z})$ by Proposition \ref{prop:finite-field-square-classes-and-universality}(iv).\\
Since $\mathcal{P}\cap\mathcal{P}^{\perp}=\{\mathbf{0}\}$ and $\mathcal{P}\perp\mathcal{P}^{\perp}$, $\mathbf{p}+\mathbf{z}\neq\mathbf{0}$ and $q(\mathbf{p}+\mathbf{z})=0$.
\item Induct on $\dim\mathcal{V}$. If $\mathcal{V}$ is anisotropic, then $\dim\mathcal{V}\leq2$ by (i); take $s=0$ and $\mathcal{R}=\mathcal{V}$.\\
Otherwise choose $\mathbf{0}\neq\mathbf{x}\in\mathcal{V}$ with $q(\mathbf{x})=0$. By Theorem \ref{thm:isotropic-vector-sits-in-hyperbolic-plane}, $\mathbf{x}$ lies in a hyperbolic plane $\mathcal{H}$.\\
By Theorem \ref{thm:metric-orthogonal-complement-dimension}(ii),(iv), $\mathcal{V}=\mathcal{H}\odot\mathcal{H}^{\perp}$, where $\mathcal{H}^{\perp}$ is nonsingular and has dimension $\dim\mathcal{V}-2$.\\
By the induction hypothesis, $\mathcal{H}^{\perp}\cong\mathcal{H}^{\odot s}\odot\mathcal{R}$ with $\mathcal{R}$ nonsingular anisotropic and $\dim\mathcal{R}\leq2$. Hence $\mathcal{V}\cong\mathcal{H}^{\odot(s+1)}\odot\mathcal{R}$.
\end{enumerate}
\hlt{Structural Commentary}: Cancel a represented value against an orthogonal line, then split the resulting hyperbolic plane inductively.\\
\hlt{Key Insight}: A nonsingular plane represents the negative norm needed to force isotropy.
\end{proof}
\end{lemma}

\begin{theorem} \label{thm:odd-finite-field-witt-decomposition} \hlt{Odd Finite-Field Orthogonal Decomposition}\\
Let $\mathbb{F}_q$ have odd characteristic, fix a nonsquare $d\in\mathbb{F}_q^{\times}$, and let $\mathcal{V}\neq\{\mathbf{0}\}$ be nonsingular orthogonal.\\
Put $S=(\mathbb{F}_q^{\times})^2$, write $\mathcal{H}^{\odot s}$ for $s$ hyperbolic planes, and let $\langle a\rangle$ denote a line of value $a$. Set $\mathcal{A}_d\defeq\langle1\rangle\odot\langle-d\rangle$.
\begin{enumerate}[label=(\roman*)]
\setlength{\itemsep}{0pt}
\item If $\dim\mathcal{V}=2m+1$, then $(\exists\,a\in\mathbb{F}_q^{\times})\,\mathcal{V}\cong\mathcal{H}^{\odot m}\odot\langle a\rangle\ \wedge\ aS=(-1)^m\Delta(\mathcal{V})$.
\item If $\dim\mathcal{V}=2m$, then $m\geq1$ and exactly one of $\mathcal{V}\cong\mathcal{H}^{\odot m}$ and $\mathcal{V}\cong\mathcal{H}^{\odot(m-1)}\odot\mathcal{A}_d$ holds.
\item If $\dim\mathcal{V}=2m$, then $\mathcal{V}\cong\mathcal{H}^{\odot m}\Longleftrightarrow\Delta(\mathcal{V})=(-1)^mS$.
\item If $\dim\mathcal{V}=2m$, then $\mathcal{V}\cong\mathcal{H}^{\odot(m-1)}\odot\mathcal{A}_d\Longleftrightarrow\Delta(\mathcal{V})=(-1)^mdS$.
\end{enumerate}
\begin{proof}
\begin{enumerate}[label=(\roman*)]
\setlength{\itemsep}{0pt}
\item[$(i)$] By Lemma \ref{lem:odd-finite-field-anisotropic-residue-bound}(ii), write $\mathcal{V}\cong\mathcal{H}^{\odot s}\odot\mathcal{R}$ with $\mathcal{R}$ anisotropic and $\dim\mathcal{R}\leq2$.\\
Since $2m+1=2s+\dim\mathcal{R}$, parity implies $\dim\mathcal{R}=1$ and $s=m$. Choose $a\in\mathbb{F}_q^{\times}$ with $\mathcal{R}\cong\langle a\rangle$.\\
Since $\Delta(\mathcal{H})=-S$, $\Delta(\mathcal{V})=(-1)^maS$, and hence $aS=(-1)^m\Delta(\mathcal{V})$.
\item[$(ii)$] Since $\mathcal{V}\neq\{\mathbf{0}\}$, $m\geq1$. By Lemma \ref{lem:odd-finite-field-anisotropic-residue-bound}(ii), write $\mathcal{V}\cong\mathcal{H}^{\odot s}\odot\mathcal{R}$ with $\dim\mathcal{R}\leq2$.\\
Since $2m=2s+\dim\mathcal{R}$, $\dim\mathcal{R}\in\{0,2\}$. If $\dim\mathcal{R}=0$, then $s=m$ and $\mathcal{V}\cong\mathcal{H}^{\odot m}$.\\
Otherwise $s=m-1$. Both $\mathcal{R}$ and $\mathcal{A}_d$ are anisotropic planes: for $\mathcal{A}_d$, $x^2-dy^2=0$ forces $x=y=0$ because $d$ is nonsquare.\\
Both occupy the unique nonhyperbolic rank-$2$ class by Theorem \ref{thm:orthogonal-canonical-forms-common-fields}(iv) and Proposition \ref{prop:hyperbolic-plane-orthogonal-criteria}.\\
Hence $\mathcal{R}\cong\mathcal{A}_d$, so at least one alternative holds. Their discriminants are $(-1)^mS$ and $(-1)^mdS$.\\
Since $d$ is nonsquare, these classes are distinct, so the alternatives are mutually exclusive.
\item[$(iii)$] ($\Rightarrow$) If $\mathcal{V}\cong\mathcal{H}^{\odot m}$, then $\Delta(\mathcal{V})=(-1)^mS$ by $\Delta(\mathcal{H})=-S$.
\item[$(iii)$] ($\Leftarrow$) By (ii), one alternative holds. Its discriminant is $(-1)^mS$ only in the first case, so $\mathcal{V}\cong\mathcal{H}^{\odot m}$.
\item[$(iv)$] ($\Rightarrow$) If $\mathcal{V}\cong\mathcal{H}^{\odot(m-1)}\odot\mathcal{A}_d$, then $\Delta(\mathcal{V})=(-1)^mdS$ by $\Delta(\mathcal{A}_d)=-dS$.
\item[$(iv)$] ($\Leftarrow$) By (ii), one alternative holds. Its discriminant is $(-1)^mdS$ only in the second case, so $\mathcal{V}\cong\mathcal{H}^{\odot(m-1)}\odot\mathcal{A}_d$.
\end{enumerate}
\hlt{Structural Commentary}: Split the anisotropic residue, determine its dimension by parity, then identify its class by discriminant.\\
\hlt{Key Insight}: Dimension parity fixes residue size; discriminant fixes its square class.
\end{proof}
\end{theorem}
